% ============================================================
% Section 6: Conclusion
% ============================================================
\section{Conclusion}
\label{sec:conclusion}

This paper presents the first systematic characterization of the data
requirements for ML-based modeling attacks on XOR Arbiter PUFs.  Our
experiments, spanning approximately 30 GPU-hours and covering XOR
complexities from 2 to 10, establish four principal findings.

First, attack success is governed by a probabilistic S-curve
$P(\text{success} \mid N, k)$ rather than a deterministic threshold.  For each
XOR level $k$, there exists a transition region in the CRP budget $N$ where
identical experimental configurations can produce either random guessing
($\approx 50\%$) or near-perfect classification ($> 99\%$).  For 7-XOR,
attack success reaches 80\% at 1.25M CRPs (8/10 trials) and 100\% at 2M
CRPs (5/5 trials).  The transition
midpoint $N_{50}$ scales by a factor of roughly 3--6$\times$ per additional
XOR level for $k = 6\to7\to8$ (the only transitions with reliable $N_{50}$
estimates), while the $k=2\to4$ transition is anomalous ($N_{50}$
decreases).  Across the full range, $N_{50}$ grows from $\sim 1 \times
10^5$ CRPs for 2-XOR to $\sim 8 \times 10^6$ CRPs for 8-XOR, with a
preliminary lower bound of $>1.5\times10^7$ CRPs for 9-XOR.  This scaling is
exponential but not super-exponential, though the per-level factor shows
variability across levels.

Second, model architecture has negligible impact on attack success below the
data threshold.  Eight distinct MLP architectures---varying in depth, width,
activation functions, residual connections, and parameter count by
$15\times$---and two evolution-strategy baselines (PSO, CMA-ES) with the
structurally correct delay parameterization all fail identically at
$\approx 50\%$ accuracy on 8-XOR with 2M and 5M CRPs.  The same standard MLP,
given 10M CRPs, achieves 99.71\%.  Data quantity,
not model design or optimization paradigm, is the decisive factor.
Transfer learning across XOR levels provides no benefit, confirming that each
XOR complexity defines a fundamentally distinct learning problem.

Third, logistic regression fails for all tested configurations with $k \geq 4$
within the trainable data range, producing random accuracy regardless of data
volume.  This confirms that the
XOR mixing function is fundamentally inaccessible to linear classifiers and
establishes a lower bound on the required model complexity for PUF modeling
attacks.

Fourth, Interpose PUFs are significantly more vulnerable than their security
claims suggest.  A $(4,4)$-iPUF, claimed equivalent to an 8-XOR APUF, is
attacked with 99.05\% accuracy using only 1M CRPs---while a true 6-XOR APUF
requires 500K CRPs for comparable reliability.  The split architecture of the
iPUF leaks information that, in our tested configurations, is consistent
with an effective security level of a
$(k_{\text{up}} + k_{\text{down}} - 2)$-XOR APUF.

These findings have direct implications for PUF security evaluation.  Security
claims should be probabilistic: instead of asking ``can $k$-XOR be broken?'',
the field should ask ``at what CRP budget does the attack succeed with what
probability?''  An 8-XOR APUF with CRP exposure limited to $5\times10^6$
provides practical security against single-GPU ML attacks, while the 9-XOR
$N_{50}$ exceeds $1.5\times10^7$ CRPs (very rough lower bound; 0/3 at 15M;
29 total trials across 5 CRP levels), beyond the memory
capacity of our 12~GB GPU: the 12M and 15M measurements were feasible only on
a 24~GB RTX~4090.  At $k = 10$, no successes were observed in three trials at
$10^7$ CRPs, and 15{,}000{,}000~CRPs fail with out-of-memory errors even on
a 24~GB GPU, establishing a practical upper bound for single-GPU attacks on
contemporary hardware.

Several limitations suggest directions for future work.  Our experiments use
simulated noise-free PUFs; silicon noise, environmental variation, and
manufacturing gradients may alter the transition characteristics.  We test
only MLP architectures and two evolution-strategy baselines; transformers,
convolutional networks, or kernel methods could exhibit different scaling
behavior.  We use uniformly random
challenges; adaptive or chosen-challenge strategies may reduce the data
requirement.  Nevertheless, our central result---that attack success is a
probabilistic function of CRP budget, not a deterministic property of
architecture or algorithm---provides a new framework for understanding and
evaluating the security of strong PUFs against ML-based modeling attacks.
